\chapter{Wprowadzenie}
% =======================================================
\epigraph{\textit{„Lotnictwo to jedyna branża, w której błąd w równaniu matematycznym może natychmiast uderzyć w ziemię z prędkością 100 kilometrów na godzinę.”}}{\ldots z inżynierskiego folkloru}

\section{Od maszyny zdalnie sterowanej do agenta autonomicznego}

Kiedy myślimy o bezzałogowych statkach powietrznych (BSP), potocznie nazywanych dronami, najczęściej wyobrażamy sobie człowieka stojącego na ziemi z aparaturą RC (Radio Control) w dłoniach. W takim modelu człowiek jest pilotem: jego układ nerwowy analizuje obraz, a dłonie wychylają drążki, korygując na bieżąco lot maszyny. 

Jednak współczesny przemysł, geodezja, nowoczesne rolnictwo czy logistyka wymuszają zmianę tego paradygmatu. Przy zadaniach wymagających milimetrowej precyzji, takich jak wielogodzinne skanowanie połaci lasu czy loty w zwartym roju, ludzka percepcja i czas reakcji stają się niewystarczające. Statek powietrzny przestaje być zdalnie sterowaną zabawką, a staje się \textbf{agentem autonomicznym}. 

W tym nowym ujęciu człowiek przestaje pełnić funkcję \textit{pilota}, a staje się \textit{nadzorcą}. Zamiast korygować wychylenie silników w czasie rzeczywistym, inżynier zadaje maszynie \textbf{cel zdefiniowany matematycznie}. To wbudowane w drona algorytmy muszą samodzielnie przeliczyć wektory przyspieszeń, skompensować wiatr i ostatecznie podjąć decyzję o tym, jaką moc przekazać na konkretne wirniki.


\section{Dlaczego matematyka wygrywa z joystickiem?}

Podstawowym celem tego skryptu jest wypełnienie luki między teoretyczną analizą matematyczną a surową fizyką programowania maszyn. Świat wirtualny, w którym na co dzień działają algorytmy informatyczne, jest wyidealizowany i przewidywalny. Świat fizyczny wręcz przeciwnie, podlega bezwładności, grawitacji i zakłóceniom zewnętrznym (takim jak szum sensorów czy podmuchy wiatru).

Aby dron mógł skutecznie połączyć te dwa światy, niezbędny jest odpowiedni system obliczeniowy. W trakcie niniejszego kursu udowodnimy, dlaczego solidny aparat matematyczny pozwala na tworzenie systemów niezawodnych. Zamiast manualnego operowania dronem, nauczymy się:

\begin{itemize}
    \item \textbf{Ciągłego planowania trajektorii} -- wykorzystania równań parametrycznych i analizy wektorowej, aby uniknąć konieczności zatrzymywania się statku w przestrzeni.
    \item \textbf{Optymalizacji misji} -- zastosowania geometrii obliczeniowej i teorii grafów (np. problemu komiwojażera) w celu maksymalizacji pokrycia terenu przy minimalnym zużyciu zasobów baterii.
    \item \textbf{Robotyki roju (Swarm Robotics)} -- przeniesienia koncepcji znanych z biologii (np. klucze ptaków) na modele układów dynamicznych (takich jak hybrydowy model Cuckera-Smale'a), gdzie lokalne interakcje matematyczne jednostek prowadzą do inteligentnych, emergentnych zachowań całej grupy, całkowicie bez udziału centralnego komputera sterującego.
\end{itemize}

Kurs ten został stworzony z myślą o studentach kierunków ścisłych, którzy posiadają już aparat analityczny, ale chcą zobaczyć, w jaki sposób zredukować go do linii kodu i wysłać w powietrze. Za pomocą profesjonalnego środowiska symulacyjnego (SITL) oraz protokołu wymiany danych (MAVLink), przejdziesz proces od zdefiniowania abstrakcyjnego równania na kartce, aż do przeprowadzenia stabilnego lotu wysoce zautomatyzowanego drona.